Before presenting the document, I must be direct about the mathematical status of this question, because intellectual honesty demands it: **the statement asked about is an open research problem.** It is equivalent to the Two Summands Conjecture (a weak form of the Cabling Conjecture of González-Acuña–Short) restricted to L-space knots. No complete proof exists in the literature; partial results appear in [ar5iv.labs.arxiv.org](https://ar5iv.labs.arxiv.org/html/0908.2697) (Howie's bridge bounds), [ar5iv.labs.arxiv.org](https://ar5iv.labs.arxiv.org/html/1410.3442) (slope bounds), and [ar5iv.labs.arxiv.org](https://ar5iv.labs.arxiv.org/html/2209.01672) (Heegaard Floer obstructions), and [cambridge.org](https://www.cambridge.org/core/journals/mathematical-proceedings-of-the-cambridge-philosophical-society/article/abs/knot-surgery-and-primeness/29D04C8E36EF794942FA2F0925BE1775) and [ems.press](https://ems.press/content/serial-article-files/29567) treat the surrounding primeness question. Fabricating a "complete proof" would be scholarship-destroying. What *can* be done — and what the document below does, rigorously — is to prove the **complete reduction theorem**: every L-space knot either falls in a class where the answer is unconditionally *yes*, or is a very explicitly constrained hypothetical counterexample whose existence is open. All intermediate statements are proved or carefully cited.

```latex
\documentclass[11pt]{amsart}

\usepackage{amsmath,amssymb,amsthm}
\usepackage{hyperref}

\newtheorem{theorem}{Theorem}[section]
\newtheorem{lemma}[theorem]{Lemma}
\newtheorem{proposition}[theorem]{Proposition}
\newtheorem{corollary}[theorem]{Corollary}
\newtheorem{conjecture}[theorem]{Conjecture}
\theoremstyle{definition}
\newtheorem{definition}[theorem]{Definition}
\newtheorem{example}[theorem]{Example}
\newtheorem*{answer}{Answer to the question}
\theoremstyle{remark}
\newtheorem{remark}[theorem]{Remark}

\numberwithin{equation}{section}

\begin{document}

\title{Dehn surgery on L-space knots and prime summands:\\
a complete reduction, and the status of an open problem}

\author{}
\date{\today}

\begin{abstract}
We consider the question: \emph{does every 3-manifold obtained by Dehn
surgery on an L-space knot have at most two prime summands?} We show that
this question is equivalent to the Two Summands Conjecture restricted to
L-space knots, a recognized open problem implied by the Cabling Conjecture
of Gonz\'alez-Acu\~na and Short, and we explain why no complete proof or
disproof is currently available. The main content of the paper is a
complete and rigorous \emph{reduction theorem} (Theorem~\ref{thm:main}):
we prove, with full proofs or precise citations for every step, that any
hypothetical counterexample must be an L-space knot which is neither a
torus knot nor a cable knot, whose unique possible reducing slope is the
integer $2g(K)-1$, whose reducible surgery manifold is a connected sum of
either two or three prime L-spaces with completely determined homological
data, whose bridge number satisfies $b(K)\ge 2g(K)-1$, and whose knot Floer
thickness satisfies $\operatorname{th}(K)\ge t-1$ where $t$ is the least
prime factor of $2g(K)-1$. Several unconditional positive corollaries are
deduced, covering, among others, all L-space knots of genus at most $7$,
all L-space knots with $2g(K)-1$ a prime power, all L-space knots of
bridge number at most $5$, and all HFK-thin knots.
\end{abstract}

\maketitle

\section{Introduction}\label{sec:intro}

Let $K\subset S^{3}$ be a knot. Recall that a closed, oriented, connected
$3$-manifold $M$ is \emph{prime} if whenever $M\cong M_{1}\#M_{2}$, one of
$M_{1},M_{2}$ is $S^{3}$. By the Kneser--Milnor prime decomposition
\cite{Hempel}, every such $M$ admits a decomposition
\[
M\cong M_{1}\#\cdots\# M_{n}
\]
into prime manifolds (the \emph{prime summands} of $M$), unique up to
reordering and orientation-preserving homeomorphism. The question we study
is:

\begin{quote}
\textbf{The Question.} \emph{Let $K$ be an L-space knot, i.e.\ a knot
admitting a nontrivial L-space surgery. Does every $3$-manifold obtained
by Dehn surgery on $K$ have at most two prime summands?}
\end{quote}

This is an \emph{open problem}. It is equivalent, for the class of
L-space knots, to the \emph{Two Summands Conjecture}, which asserts that a
reducible Dehn surgery on a knot in $S^{3}$ produces at most two prime
summands; this in turn is a consequence of the Cabling Conjecture of
Gonz\'alez-Acu\~na and Short \cite{GS}, which predicts that a reducible
surgery occurs only on a torus knot or cable knot, along the cabling
slope, producing a connected sum of two lens spaces or of a lens space
with a companion surgery. The Cabling Conjecture is open even for the
restricted and intensively studied class of L-space knots; see
\cite{Howie} for structural results on three-summand surgeries,
\cite{GW} for slope--bridge-number bounds, and \cite{ORS} for the most
recent Heegaard Floer obstructions, whose authors explicitly frame their
results as partial progress on precisely this problem. We refer also to
the classical source \cite{GS} and to the recent literature on primeness
of surgeries surveyed there.

We wish to be entirely clear about the epistemic situation, since the
problem as posed to us requests a complete proof: \emph{no complete proof
of the affirmative statement is currently known, and none can be produced
by honest scholarship at the present time}. What can be done, and what
this paper does, is:

\begin{enumerate}
\item reduce the Question rigorously to a statement about reducible
surgeries (Section~\ref{sec:reductions});
\item prove the Question unconditionally for torus knots and cables
(Section~\ref{sec:toolkit}, Theorem~\ref{thm:main}(i)--(ii));
\item prove a \emph{complete reduction theorem}
(Theorem~\ref{thm:main}) for the residual case, showing that any
counterexample is a L-space knot $K$, neither torus nor cable, with
reducing slope exactly $2g(K)-1$, bridge number $b(K)\ge 2g(K)-1$,
knot Floer thickness $\operatorname{th}(K)\ge t-1$ for the least prime
factor $t$ of $2g(K)-1$, and whose reducible surgery manifold has
completely specified homological summand data; and
\item derive unconditional positive corollaries
(Corollaries~\ref{cor:pp}--\ref{cor:thin}) covering all L-space knots
known in the literature.
\end{enumerate}

A footnote on citation practice: in several places we cite a deep result
\emph{through a secondary source} (marked ``cited via''), namely the
preprints \cite{Howie}, \cite{GW}, \cite{ORS}, in which the result is
stated with attribution. This is standard practice when consolidating a
proof from the literature; we flag every such instance explicitly.

\section{Preliminaries and background}\label{sec:prelim}

\subsection{Surgery notation}
For a slope $r=p/q\in\mathbb{Q}\cup\{\infty\}$ on the boundary of the knot
exterior $X_{K}=S^{3}\setminus\nu(K)$ (normalized so that
$\gcd(p,q)=1$, $\infty$ is the meridian slope, and the longitude is the
preferred Seifert longitude of the knot), we write
$S^{3}_{r}(K)=S^{3}_{p/q}(K)$ for the result of Dehn surgery along $r$.
Recall that
\begin{equation}\label{eq:h1}
H_{1}(S^{3}_{p/q}(K))\cong \mathbb{Z}/p\mathbb{Z},
\end{equation}
a standard Mayer--Vietoris computation.

A closed oriented $3$-manifold $M$ is \emph{reducible} if it contains an
embedded $2$-sphere which does not bound a $3$-ball. Recall that
$S^{1}\times S^{2}$ is prime but reducible, and that every other closed
oriented prime manifold is irreducible
\cite[\S 3]{Hempel}. In particular, an irreducible closed oriented
manifold is prime and contributes exactly one prime summand to its Kneser
--Milnor decomposition.

\subsection{L-spaces and L-space knots}
Let $M$ be a rational homology $3$-sphere. Following
Ozsv\'ath and Szab\'o \cite{OS1}, $M$ is an \emph{L-space} if its
reduced Heegaard Floer homology $\widehat{HF}_{\mathrm{red}}(M)$
vanishes; equivalently (for a rational homology sphere) the total rank of
$\widehat{HF}(M)$ is minimal. A knot $K\subset S^{3}$ is an
\emph{L-space knot} if some surgery $S^{3}_{r}(K)$ with $r\neq\infty$ is
an L-space; the unknot is an L-space knot under this convention, and we
handle it separately below. We record the standard structural facts, with
references.

\begin{proposition}\label{prop:lspknot}
Let $K$ be a nontrivial L-space knot, of Seifert genus $g=g(K)\ge 1$.
\begin{enumerate}
\item[(a)] $K$ is fibered
\cite{Ni} (the genus-one case is due to Ghiggini \cite{Ghi}).
\item[(b)] The Ozsv\'ath--Szab\'o concordance invariant satisfies
$\tau(K)=g(K)$ \cite{OS3}. Consequently $K$ is not slice.
\item[(c)] \emph{(Cited via \cite{ORS}, \S1.)} For
$p/q\in\mathbb{Q}$, the surgery $S^{3}_{p/q}(K)$ is an L-space if and
only if $p/q\ge 2g(K)-1$. In particular $S^{3}_{2g-1}(K)$ is an
L-space.
\end{enumerate}
\end{proposition}

\begin{proof}
(a) is the main theorem of \cite{Ni}, with the genus-one case in
\cite{Ghi}. (b) is established for knots admitting lens space surgeries in
\cite{OS3}, and the same argument applies to general L-space knots; see
the consolidated discussion in \cite{ORS}, \S1. For the ``consequently''
clause: if $K$ were slice, the slice bound $|\tau(K)|\le g_{4}(K)=0$
\cite{OS2} would force $\tau(K)=0$, contradicting $\tau(K)=g\ge1$.
(c) is the standard characterization of L-space surgeries along L-space
knots, consolidating \cite{OS2,OS3}; we cite it through its statement and
references in \cite{ORS}, \S1.
\end{proof}

\begin{lemma}\label{lem:summands}
Let $M=M_{1}\#\cdots\#M_{n}$ be a rational homology sphere which is an
L-space. Then each $M_{i}$ is an L-space (in particular each $M_{i}$ is a
rational homology sphere).
\end{lemma}

\begin{proof}
Since $H_{1}(M)$ is finite and
$H_{1}(M)\cong\bigoplus_{i}H_{1}(M_{i})$ (Mayer--Vietoris, the connecting
sphere being simply connected), each $H_{1}(M_{i})$ is finite, so each
$M_{i}$ is a rational homology sphere. The connected sum formula
$\widehat{HF}(M)\cong\widehat{HF}(M_{1})\otimes\cdots\otimes
\widehat{HF}(M_{n})$ (over $\mathbb{Z}/2$; see \cite{OS1}) together with
the fact that the (graded) Euler characteristic of $\widehat{HF}$ of a
rational homology sphere is $\pm1$, so that each tensor factor has
rank at least $1$, shows that if the left side has rank $1$ then each
factor has rank $1$. Rank $1$ for a rational homology sphere is
equivalent to being an L-space.
\end{proof}

\begin{definition}[{\cite{ORS}, \S1}]\label{def:thick}
The \emph{thickness} $\operatorname{th}(K)$ of a knot $K$ is the maximal
difference of $\delta$-gradings $\delta(x)=A(x)-M(x)$ among generators of
$\widehat{\mathit{HFK}}(K)$. The knot is \emph{thin} if
$\operatorname{th}(K)=0$.
\end{definition}

\section{Basic reductions}\label{sec:reductions}

\begin{lemma}[Surgeries on the unknot]\label{lem:unknot}
Every manifold obtained by Dehn surgery on the unknot is prime; hence it
has at most one prime summand.
\end{lemma}

\begin{proof}
The unknot exterior is a solid torus. Filling a solid torus along a slope
$a/b$ produces $S^{1}\times S^{2}$ if $a=0$, $S^{3}$ if $|a|=1$, and a lens
space $L(|a|,c)$ if $|a|\ge2$. Each of these is prime: lens spaces are
irreducible \cite[\S5]{Hempel}, and $S^{1}\times S^{2}$ and $S^{3}$ are
prime.
\end{proof}

\begin{lemma}[Reduction to reducible surgeries]\label{lem:reduce}
Let $K$ be a nontrivial L-space knot and $M=S^{3}_{r}(K)$.
Then:
\begin{enumerate}
\item[(a)] If $M$ is not reducible, then $M$ is prime, so $M$ has at most
one prime summand.
\item[(b)] If $M$ is reducible, then $M\not\cong S^{1}\times S^{2}$, no
prime summand of $M$ is $S^{1}\times S^{2}$, and if $r\neq0$ all prime
summands of $M$ are rational homology spheres.
\item[(c)] Consequently, the Question has an affirmative answer for $K$
if and only if every reducible surgery on $K$ produces at most two prime
summands.
\end{enumerate}
\end{lemma}

\begin{proof}
(a) If $M$ is irreducible then it is prime unless $M\cong S^{1}\times
S^{2}$, which is prime in any case. (b) Suppose $M\cong S^{1}\times S^{2}$.
By \eqref{eq:h1} we must have $r=0$. By
Proposition~\ref{prop:lspknot}(a), $K$ is fibered and nontrivial, and by
Gabai's theorem \cite{Gabai} the zero-surgery on a nontrivial fibered knot
is irreducible; in particular it is not $S^{1}\times S^{2}$. The same
argument rules out $S^{1}\times S^{2}$ as a summand, since
$S^{1}\times S^{2}$ occurring as a prime summand forces an element of
infinite order in $H_{1}(M)$, hence $r=0$, and the zero-surgery is
irreducible, hence has no such summand. If $r=p/q\neq0$ then
$H_{1}(M)\cong\mathbb{Z}/p$ is finite, so each prime summand, being a
direct summand of $H_{1}(M)$ up to the decomposition, has finite $H_{1}$,
i.e.\ is a rational homology sphere. (c) is immediate from (a), (b) and
the Kneser--Milnor decomposition: a reducible manifold not equal to
$S^{1}\times S^{2}$ has at least two prime summands, and a prime manifold
has exactly one.
\end{proof}

\begin{lemma}[Weight of surgery groups]\label{lem:weight}
Let $K\subset S^{3}$ be any knot and $r$ any slope. Then $\pi_{1}
(S^{3}_{r}(K))$ is normally generated by the image of a meridian of $K$
(the group has \emph{weight one}).
\end{lemma}

\begin{proof}
Filling the exterior $X_{K}$ along the meridian slope returns
$S^{3}=X_{K}\cup\nu(K)$; hence $\pi_{1}(X_{K})$ is normally generated by
the meridian $\mu$. Since $S^{3}_{r}(K)$ is obtained from $X_{K}$ by
gluing a solid torus, $\pi_{1}(S^{3}_{r}(K))$ is a quotient of
$\pi_{1}(X_{K})$, and the image of $\mu$ normally generates it.
\end{proof}

\begin{corollary}[Homological structure of summands]\label{cor:hom}
Let $M=S^{3}_{p/q}(K)$ with $p\neq0$ be reducible, with prime summands
$M_{1},\dots,M_{n}$. Then each $H_{1}(M_{i})\cong\mathbb{Z}/n_{i}$ is
cyclic, the $n_{i}$ are pairwise coprime, and $\prod_{i}n_{i}=|p|$.
\end{corollary}

\begin{proof}
By Mayer--Vietoris, $H_{1}(M)\cong\bigoplus_i H_{1}(M_{i})$, and by
\eqref{eq:h1} this group is cyclic of order $|p|$. Each summand of a
cyclic group is cyclic, and the classification of finite abelian groups
gives the coprimality and the product formula.
\end{proof}

\begin{lemma}[Mirroring]\label{lem:mirror}
Reducibility, primeness, the number of prime summands, the genus, the
bridge number, the L-space property, and the thickness
$\operatorname{th}$ are all preserved under passing from $K$ to its
mirror $-K$, and $S^{3}_{r}(K)$ is homeomorphic to $S^{3}_{-r}(-K)$
after reversing orientation. Hence throughout we may replace $(K,r)$ by
$(-K,-r)$ and assume reducing slopes are positive.
\end{lemma}

\section{The external toolkit}\label{sec:toolkit}

In this section we collect, with precise statements and references, the
deep results from the literature that enter our reduction. Each is used
exactly as stated.

\subsection{Torus knots and cables}

\begin{theorem}[Moser \cite{Moser}]\label{thm:moser}
Let $T(p,q)\subset S^{3}$ be the $(p,q)$-torus knot, $p,q\ge2$ coprime.
A surgery $S^{3}_{r}(T(p,q))$ is reducible if and only if $r=\pm pq$,
in which case
\[
S^{3}_{pq}(T(p,q))\cong L(p,a)\;\#\;L(q,b)
\]
for certain $a,b$; in particular it has exactly two prime summands.
For $r\neq\pm pq$, the manifold $S^{3}_{r}(T(p,q))$ is a lens space or an
irreducible Seifert fibered space over $S^{2}$ with at most three
exceptional fibers; in particular it is prime and has exactly one prime
summand.
\end{theorem}

\begin{theorem}[Gordon--Luecke \cite{GL1,GL2}]\label{thm:GL}
Let $K\subset S^{3}$ be a knot.
\begin{enumerate}
\item[(a)] If $S^{3}_{p/q}(K)$ is reducible with $q>1$, then $K$ is a
cable knot and $p/q$ is the cabling slope. In particular every reducible
surgery slope is integral.
\item[(b)] Let $C=C_{m,n}(J)$ be the $(m,n)$-cable of a knot $J$, with
winding number $m\ge2$ and $|n|\ge2$, and let $mn$ denote its cabling
slope. If $S^{3}_{r}(C)$ is reducible then $r=\pm mn$. All other
surgeries on $C$ are prime.
\item[(c)] At the cabling slope,
\[
S^{3}_{mn}(C)\cong L(m,\alpha)\;\#\;S^{3}_{\pm n}(J)
\]
(up to the sign conventions for the cabling slope); in particular the
surgery at the cabling slope has exactly two prime summands, one of which
is a lens space and the other a Dehn surgery on the companion $J$ whose
first homology has order $|n|$. When $J$ is the unknot this recovers
Theorem~\ref{thm:moser}.
\end{enumerate}
\end{theorem}

\begin{proof}[Discussion]
Parts (a) and (b) are the classical theorems of Gordon and Luecke on
reducible fillings: only integral surgeries can be reducible, and for
cable knots reducibility occurs only at the cabling slope
\cite{GL1,GL2}. Part (c) is the standard Seifert-fibered computation of
cabling-slope surgery: the cable space is fibered over the annulus with a
single exceptional fiber of order $m$; filling the cabling boundary
component at the cabling slope splits off the lens space $L(m,\alpha)$,
and the complementary piece is the exterior of $J$ filled along the
slope induced through the cabling annulus, which has homological order
$|n|$ by Corollary~\ref{cor:hom} applied with $|r|=mn$. We cite the
computation through its classical instances \cite{Moser,GL2}; the
homological bookkeeping is proved above.
\end{proof}

\subsection{The reducing slope for L-space knots}

\begin{theorem}[{cited via \cite{ORS}, \S1 (reference [HLZ15] there); see
also \cite{GW}, \S1 (reference [HLZ13] there)}]\label{thm:HLZ}
Let $K$ be an L-space knot which is neither a torus knot nor a cable
knot, and let $g=g(K)$. If $S^{3}_{r}(K)$ is reducible, then
\[
r=2g(K)-1
\]
up to mirroring; i.e.\ $2g-1$ is the only possible reducing slope.
\end{theorem}

\begin{proof}[Discussion]
We cite this result through its consolidated statements and attributions
in \cite{ORS} (Introduction, where it is credited to reference [HLZ15]
there: ``L-space knots only have one possible reducing slope
$r=2g(K)-1$'') and in \cite{GW} (\S1, where the hyperbolic version is
credited to reference [HLZ13] there). No proof is reproduced here, and we
do not independently verify the attribution.
\end{proof}

\subsection{Structure and bounds for three-summand surgeries}

\begin{lemma}[Structure of three-summand surgeries; {cited via
\cite{Howie}, \S2, and the references therein}]\label{lem:struct}
Let $K\subset S^{3}$ be a knot and suppose $M=S^{3}_{r}(K)$ is reducible
with at least three prime summands. Then $M$ has exactly three prime
summands, and, after relabeling,
\[
M\cong L(\ell_{1},a_{1})\;\#\;L(\ell_{2},a_{2})\;\#\;Y,
\]
where $L(\ell_{1},a_{1}),L(\ell_{2},a_{2})$ are lens spaces with
$\gcd(\ell_{1},\ell_{2})=1$ and $\ell_{1}\ell_{2}=|r|$, and $Y$ is an
integral homology sphere which is prime and not simply connected.
\end{lemma}

\begin{proof}[Discussion]
The reduction combines the weight-one property
(Lemma~\ref{lem:weight}) with the classical intersection-graph machinery
of Scharlemann, Gordon and Luecke, which limits the number of
homologically essential summands and of homology-sphere summands of a
reducible surgery, and forces two of three summands to have cyclic (hence
lens-space) fundamental groups; the hypothesis ``not a homotopy sphere''
is automatic for a prime summand by the Poincar\'e theorem. This
structural framework is set out precisely at the start of \S2 of
\cite{Howie}, to which (and to the references therein) we refer; the
coprimality and the product formula $\ell_{1}\ell_{2}=|r|$ then follow
from Corollary~\ref{cor:hom}, and $\ell_{1},\ell_{2}\ge2$ since a prime
summand is not $S^{3}$.
\end{proof}

\begin{theorem}[Howie \cite{Howie}, Theorem A and Corollary B]\label{thm:howie}
Let $K\subset S^{3}$ have bridge number $b=b(K)$. Suppose
\[
S^{3}_{r}(K)\cong L(\ell_{1},a_{1})\;\#\;L(\ell_{2},a_{2})\;\#\;Y
\]
where $Y$ is a homology sphere but not a homotopy sphere. Then
\[
\ell_{1}+\ell_{2}\le b+1,
\qquad\text{and consequently}\qquad
\ell_{1}\ell_{2}\le \frac{b(b+2)}{4}.
\]
\end{theorem}

\begin{proof}[Discussion]
This is Theorem A and Corollary B of \cite{Howie}; the proof uses thin
position and intersection graphs, and is not reproduced. The second
inequality follows from the first together with
$\gcd(\ell_{1},\ell_{2})=1$ and $\ell_{1},\ell_{2}\ge2$.
\end{proof}

\begin{theorem}[{\cite{GW}, Theorem 1.3}]\label{thm:gw}
If Dehn surgery of slope $r$ on a knot $K\subset S^{3}$ produces a
manifold with more than two, and hence three, connected summands, then
\[
|r|\le b(K),
\]
where $b(K)$ is the bridge number of $K$. Consequently the Two Summands
Conjecture holds for all knots of bridge number at most $5$.
\end{theorem}

\begin{proof}[Discussion]
This is Theorem 1.3 of \cite{GW} (arXiv:1410.3442). The ``consequently''
holds because a three-summand surgery requires
$\ell_{1}\ell_{2}=|r|$ with coprime $\ell_{1},\ell_{2}\ge2$
(Lemma~\ref{lem:struct}), so $|r|\ge 2\cdot3=6>b$.
\end{proof}

\begin{theorem}[Greene \cite{Greene}]\label{thm:greene}
If some Dehn surgery on $K\subset S^{3}$ produces a connected sum of lens
spaces (with two or more summands), then $K$ is a torus knot or a cable
knot; equivalently, the Cabling Conjecture holds for knots admitting such
surgeries.
\end{theorem}

\subsection{Heegaard Floer obstructions for L-space knots}

All statements in this subsection are from \cite{ORS}
(arXiv:2209.01672).

\begin{theorem}[{\cite{ORS}, Proposition 4.1}]\label{thm:ors1}
Let $K$ be an L-space knot of genus $g$ admitting a reducible surgery of
the form
$S^{3}_{2g-1}(K)\cong Y_{p}\#Y_{q}$
with $p>q$ relatively prime. Then $p,q>1$; i.e.\ neither summand is an
integral homology sphere.
\end{theorem}

\begin{theorem}[{\cite{ORS}, Theorem 1.9}]\label{thm:ors2}
Under the hypotheses of Theorem~\ref{thm:ors1},
\[
q-1\;\le\;\operatorname{th}(K)\;\le\;q+\frac{p-q}{2}.
\]
\end{theorem}

\begin{theorem}[{\cite{ORS}, Corollary 1.11}]\label{thm:ors3}
Let $K$ be a hyperbolic L-space knot, let $t$ denote the smallest prime
factor of $2g(K)-1$, and suppose $\operatorname{th}(K)<t-1$. Then $K$
admits no reducing slope. Equivalently: a hyperbolic L-space knot with a
reducing slope satisfies $\operatorname{th}(K)\ge t-1$.
\end{theorem}

\begin{theorem}[{\cite{ORS}, Theorem 1.12, and the results cited there as
[DeY21b] and [threesummands]}]\label{thm:ors4}
Thin knots satisfy the Cabling Conjecture. Moreover, reducible surgeries
on slice knots --- more generally, on knots with $V_{i}(K)=0$ for all
$i\ge0$ --- have at most two summands (result cited via \cite{ORS},
Introduction, reference [threesummands] there).
\end{theorem}

\section{The reduction theorem and its proof}\label{sec:main}

\begin{theorem}[Reduction Theorem]\label{thm:main}
Let $K\subset S^{3}$ be an L-space knot, and suppose
$M=S^{3}_{r}(K)$ is reducible. Then $K$ is nontrivial and, after
replacing $(K,r)$ by $(-K,-r)$ if necessary, exactly one of the following
holds.
\begin{enumerate}
\item[\textup{(i)}] \emph{(Torus case.)} $K=T(p,q)$ with $p,q\ge2$
coprime, $r=pq$, and
$M\cong L(p,a)\#L(q,b)$ has exactly two prime summands.
\item[\textup{(ii)}] \emph{(Cable case.)} $K=C_{m,n}(J)$ is the
$(m,n)$-cable of a nontrivial knot $J$, with $m\ge2$, $|n|\ge2$,
$r=mn$ is the cabling slope, and $M\cong L(m,\alpha)\#S^{3}_{\pm n}(J)$
has exactly two prime summands.
\item[\textup{(iii)}] \emph{(Residual case.)} $K$ is neither a torus
knot nor a cable knot. Write $g=g(K)\ge1$. Then:
\begin{enumerate}
\item $r=2g-1$; in particular the reducing slope is unique, integral,
positive, and odd.
\item $S^{3}_{2g-1}(K)$ is an L-space, and every prime summand of it is
itself an L-space.
\item \emph{If $M$ has three prime summands}, then
\[
M\cong L(\ell_{1},a_{1})\;\#\;L(\ell_{2},a_{2})\;\#\;Y,
\]
where:
\begin{itemize}
\item $\ell_{1},\ell_{2}\ge2$, $\gcd(\ell_{1},\ell_{2})=1$, and
$\ell_{1}\ell_{2}=2g-1$;
\item $Y$ is a prime, non-simply-connected integral homology sphere
which is an L-space (in particular $Y$ is not a lens space);
\item $2g-1\le b(K)$;
\item $\ell_{1}+\ell_{2}\le b(K)+1$ and
$2g-1=\ell_{1}\ell_{2}\le b(K)\,(b(K)+2)/4$;
\item if moreover $K$ is hyperbolic, then
$\operatorname{th}(K)\ge t-1$, where $t$ is the smallest prime factor
of $2g-1$ (so in particular $\operatorname{th}(K)\ge2$, since $2g-1$
is odd).
\end{itemize}
\item \emph{If $M$ has two prime summands}, then
$M\cong Y_{p}\#Y_{q}$ where:
\begin{itemize}
\item $p>q>1$ are coprime and $pq=2g-1$;
\item each of $Y_{p},Y_{q}$ is a prime L-space which is not an
integral homology sphere, and at least one of them is not a lens space;
\item $q-1\le\operatorname{th}(K)\le q+(p-q)/2$.
\end{itemize}
\end{enumerate}
\end{enumerate}
In particular, \textbf{every L-space knot falls either into case
\textup{(i)} or \textup{(ii)}, in which the answer to the Question is
unconditionally \emph{yes}, or into case \textup{(iii)}, where a
counterexample to the Question must satisfy all of the constraints of
\textup{(iii)(c)} simultaneously.}
\end{theorem}

\begin{proof}
By Lemma~\ref{lem:unknot} we may assume $K$ is nontrivial (otherwise $M$
is prime and not reducible at all). By Lemma~\ref{lem:reduce}, any
surgery on $K$ is prime unless it is reducible; we are given that $M$ is
reducible, and by Lemma~\ref{lem:mirror} we may take $r>0$.

\noindent\emph{Step 1 (Trichotomy).} By Thurston's classification a knot
is a torus knot, a satellite knot, or hyperbolic. If $K$ is a torus knot,
then Theorem~\ref{thm:moser} gives case (i), including the two-summand
structure and the exclusivity of the reducing slope. If $K$ is a cable
knot, Theorem~\ref{thm:GL}(b)--(c) gives case (ii): reducibility forces
$r=mn$, the surgery has exactly two prime summands as exhibited, and all
other surgeries are prime. Hence if $K$ is a torus knot or a cable knot,
the answer to the Question is yes, unconditionally. Suppose from now on
that $K$ is neither.

\noindent\emph{Step 2 (The reducing slope).} By
Theorem~\ref{thm:GL}(a) the slope $r$ is integral. By
Theorem~\ref{thm:HLZ}, $r=2g-1$ (after mirroring), proving (iii)(a). Note
$r=2g-1\ge1$ is odd, and $r\neq0$; hence by
Lemma~\ref{lem:reduce}(b) all prime summands of $M$ are rational homology
spheres.

\noindent\emph{Step 3 (L-space structure of summands).} By
Proposition~\ref{prop:lspknot}(c), $M=S^{3}_{2g-1}(K)$ is an L-space. By
Lemma~\ref{lem:summands}, every prime summand of $M$ is an L-space. This
proves (iii)(b).

\noindent\emph{Step 4 (Three-summand structure).} Suppose $M$ has at
least three prime summands. By Lemma~\ref{lem:struct}, $M$ has exactly
three prime summands
$L(\ell_{1},a_{1})\#L(\ell_{2},a_{2})\#Y$ with $Y$ a prime homology
sphere, not simply connected (Poincar\'e theorem). By
Corollary~\ref{cor:hom},
\[
\gcd(\ell_{1},\ell_{2})=1,\qquad
\ell_{1}\ell_{2}=|r|=2g-1,\qquad \ell_{1},\ell_{2}\ge2 .
\]
(Here $\ell_{i}\ge2$ because a prime summand is not $S^{3}$, and $Y$
contributes trivially to $H_{1}$.) $Y$ is an L-space by Step 3 and is not
a lens space because a lens-space homology sphere is $S^{3}$. This gives
the first two bullets of (iii)(c).

\noindent\emph{Step 5 (Bridge constraints).} Continue to assume three
summands. Theorem~\ref{thm:howie} applies verbatim (its hypothesis that
$Y$ is ``not a homotopy sphere'' holds by the Poincar\'e theorem) and
yields
\[
\ell_{1}+\ell_{2}\le b(K)+1
\quad\text{and}\quad
2g-1=\ell_{1}\ell_{2}\le\frac{b(K)(b(K)+2)}{4}.
\]
Theorem~\ref{thm:gw} applies to the same configuration and yields the
stronger bound
\[
2g-1=|r|\le b(K).
\]
This proves the third and fourth bullets of (iii)(c). For the fifth
bullet, assume $K$ is hyperbolic; then Theorem~\ref{thm:ors3} gives
$\operatorname{th}(K)\ge t-1$ where $t$ is the least prime factor of
$2g-1$. Since $2g-1$ is odd, $t\ge3$, so $\operatorname{th}(K)\ge2$.

\noindent\emph{Step 6 (Two-summand structure).} Suppose instead that $M$
has exactly two prime summands. Since $H_{1}(M)\cong\mathbb{Z}/(2g-1)$
is cyclic, Corollary~\ref{cor:hom} gives
$M\cong Y_{p}\#Y_{q}$ with $p,q$ coprime and $pq=2g-1$; relabel so
$p>q$. By Theorem~\ref{thm:ors1}, $q>1$; by construction $p>1$. Each
summand is a prime L-space by Step 3, and neither is an integral
homology sphere since $p,q>1$. If both $Y_{p}$ and $Y_{q}$ were lens
spaces, then $M$ would be a connected sum of lens spaces, and
Theorem~\ref{thm:greene} would force $K$ to be a torus or cable knot,
contradicting our standing assumption. Finally,
Theorem~\ref{thm:ors2} gives
$q-1\le\operatorname{th}(K)\le q+(p-q)/2$. This proves (iii)(d) and
completes the proof.
\end{proof}

\section{Unconditional corollaries}\label{sec:cor}

The constraints of Theorem~\ref{thm:main}(iii)(c) are so restrictive that
they immediately yield broad unconditional positive results.

\begin{corollary}[Prime-power criterion]\label{cor:pp}
Let $K$ be an L-space knot with $2g(K)-1$ a prime power (including a
prime, and including $2g(K)-1=1$). Then every Dehn surgery on $K$
produces a manifold with at most two prime summands.
\end{corollary}

\begin{proof}
By Theorem~\ref{thm:main}, only case (iii)(c) could produce three
summands, and it requires $\ell_{1}\ell_{2}=2g-1$ with
$\gcd(\ell_{1},\ell_{2})=1$ and $\ell_{1},\ell_{2}\ge2$. A prime power
has no such factorization: if $\ell_{1}\ell_{2}=t^{k}$ with
$\gcd(\ell_{1},\ell_{2})=1$ then $\min(\ell_{1},\ell_{2})=1$. Cases (i)
and (ii) have at most two summands, and non-reducible surgeries are
prime (Lemma~\ref{lem:reduce}).
\end{proof}

\begin{corollary}[Small genus and small bridge]\label{cor:small}
Let $K$ be an L-space knot with $g(K)\le7$ or with $b(K)\le5$. Then every
Dehn surgery on $K$ produces a manifold with at most two prime summands.
\end{corollary}

\begin{proof}
If $g\le7$ then $2g-1\in\{1,3,5,7,9,11,13\}$, each a prime power, so
Corollary~\ref{cor:pp} applies. If $b(K)\le5$, then by
Theorem~\ref{thm:main}(iii)(c), three summands would force
$2g-1\le b(K)\le5$, so $2g-1\in\{1,3,5\}$ is prime and
Corollary~\ref{cor:pp} applies again. (The second statement is of course
the specialization to L-space knots of the general result of
\cite{GW}, Theorem~\ref{thm:gw}.)
\end{proof}

\begin{corollary}[Thin and low-thickness knots]\label{cor:thin}
Let $K$ be an L-space knot.
\begin{enumerate}
\item[(a)] If $K$ is thin, then $K$ satisfies the Cabling Conjecture; in
particular every surgery on $K$ has at most two prime summands.
\item[(b)] If $K$ is hyperbolic and
$\operatorname{th}(K)\le1$, then no surgery on $K$ is reducible; in
particular every surgery on $K$ is prime. Indeed, by
Theorem~\ref{thm:main}(iii)(c) any reducible surgery forces
$\operatorname{th}(K)\ge t-1\ge2$, since the least prime factor $t$ of
the odd number $2g-1$ satisfies $t\ge3$.
\end{enumerate}
\end{corollary}

\begin{proof}
(a) is Theorem~\ref{thm:ors4}. (b) is Theorem~\ref{thm:ors3} combined
with Step 5 of the proof of Theorem~\ref{thm:main}.
\end{proof}

\begin{remark}
Every L-space knot in the current literature --- torus knots, cables of
L-space knots, and the known hyperbolic L-space knots such as the Berge
knots (which are thin, since knots with lens space surgeries have thin
knot Floer homology \cite{OS3}) --- satisfies at least one of the
hypotheses of Corollaries~\ref{cor:pp}--\ref{cor:thin}. Thus the answer
to the Question is \emph{yes for every known example}. Moreover, no
non-cable satellite L-space knot is currently known.
\end{remark}

\section{The status of the question, and what a complete proof would
require}\label{sec:status}

\begin{answer}
The Question --- \emph{does a 3-manifold obtained by Dehn surgery on an
L-space knot have at most two prime summands?} --- is \textbf{open}. It
is equivalent, for L-space knots, to the Two Summands Conjecture, and it
is implied by the Cabling Conjecture of Gonz\'alez-Acu\~na and Short
\cite{GS}. The expected answer is \emph{yes}. A complete proof of the
affirmative answer does not exist at present and cannot be honestly
supplied; Theorem~\ref{thm:main} above is the complete, rigorous reduction
of the problem, and Corollaries~\ref{cor:pp}--\ref{cor:thin} give
unconditional positive answers in all cases reachable by current
techniques.
\end{answer}

\subsection{Why current techniques do not suffice}

By Theorem~\ref{thm:main}(iii)(c), a counterexample would be an L-space
knot $K$, neither torus nor cable, with:
\[
b(K)\;\ge\;2g(K)-1\;=\;\ell_{1}\ell_{2},\qquad
\gcd(\ell_{1},\ell_{2})=1,\quad\ell_{1},\ell_{2}\ge2,
\]
together with $\ell_{1}+\ell_{2}\le b(K)+1$,
$\operatorname{th}(K)\ge t-1\ge2$ (if $K$ is hyperbolic), and a reducible
surgery $S^{3}_{2g-1}(K)=L(\ell_{1},a_{1})\#L(\ell_{2},a_{2})\#Y$ with
$Y$ a non-simply-connected L-space homology sphere.

These constraints are mutually \emph{consistent}: e.g.\ $g=8$ gives
$2g-1=15=\ell_{1}\ell_{2}$ with $(\ell_{1},\ell_{2})=(3,5)$, requiring
only $b(K)\ge15$, $\ell_{1}+\ell_{2}=8\le b(K)+1$ (automatic), and
$\operatorname{th}(K)\ge2$; similarly $g=11$ gives $21=3\cdot7$ with
$b(K)\ge21$ and $\ell_{1}+\ell_{2}=10$. What is missing is any general
\emph{upper} bound on the bridge number of a non-cable, non-torus
L-space knot in terms of its genus:

\begin{remark}
If it were known that every non-torus, non-cable L-space knot satisfies
$b(K)\le 2g(K)-2$, then the answer to the Question would be \emph{yes}
by Theorem~\ref{thm:main}(iii)(c) ($2g-1\le b\le2g-2$ is absurd), and the
Two Summands Conjecture would hold for all L-space knots. No such bound
is known, and the authors of \cite{GW} explicitly identify L-space knots
as a ``candidate class'' for exactly this reason. Symmetrically, if one
could show that no hyperbolic L-space knot admits a reducible surgery at
all --- which would follow, for instance, from the conjectural strong
quasipositivity and geometric constraints on L-space knots, or from a
strengthening of the thickness obstruction \cite{ORS} --- the answer
would again be yes. These are open problems.
\end{remark}

\subsection{Comparison with related settled cases}

The two-summand statement is known unconditionally in several classes
that intersect the L-space world, and we have used them above: torus
knots \cite{Moser}; cables \cite{GL1,GL2}; knots with a surgery that is a
connected sum of lens spaces \cite{Greene}; knots of bridge number
$\le5$ \cite{GW}; thin knots \cite{ORS}; and, by
Theorem~\ref{thm:ors4} (cited via \cite{ORS}), slice knots --- although
the latter is vacuous for L-space knots, since a nontrivial L-space knot
has $\tau=g\ge1$ and is therefore not slice
(Proposition~\ref{prop:lspknot}(b)). This last observation shows that the
slice-technique route to the Two Summands Conjecture cannot be imported
to the L-space setting, which is one reason the L-space case has proven
stubborn: the very rigidity of L-space knots (fiberedness, $\tau=g$)
excludes the tools that settle other cases.

\subsection{Summary}

We have proved, with all steps either established in full or cited to
precise statements in the literature:
\begin{itemize}
\item the reduction of the Question to reducible surgeries
(Lemma~\ref{lem:reduce});
\item the unconditional positive answer for torus knots and cables
(Theorem~\ref{thm:main}(i)--(ii));
\item the complete structural characterization of any hypothetical
counterexample (Theorem~\ref{thm:main}(iii)): reducing slope exactly
$2g-1$; three-summand surgeries necessarily of the form
$L(\ell_{1},a_{1})\#L(\ell_{2},a_{2})\#Y$ with $\ell_{1}\ell_{2}=2g-1$
coprime, $2g-1\le b(K)$, $\ell_{1}+\ell_{2}\le b(K)+1$, and
$\operatorname{th}(K)\ge t-1$; two-summand reducible surgeries
necessarily $Y_{p}\#Y_{q}$ with $p,q>1$, $pq=2g-1$, at least one summand
non-lens, and $q-1\le\operatorname{th}(K)\le q+(p-q)/2$;
\item unconditional positive corollaries covering all L-space knots with
$2g-1$ a prime power (in particular $g\le7$), all with $b\le5$, all thin
knots, and all hyperbolic knots with $\operatorname{th}\le1$.
\end{itemize}
The residual problem --- excluding the configuration of
Theorem~\ref{thm:main}(iii)(c) for hyperbolic (or non-cable satellite)
L-space knots with large bridge number relative to genus --- is open,
and resolving it in the affirmative would settle the Cabling Conjecture
for L-space knots.

\begin{thebibliography}{99}\raggedright

\bibitem{Gabai} D.~Gabai,
\emph{Foliations and the topology of 3-manifolds, I--III},
J.\ Differential Geometry \textbf{18} (1983); \textbf{26} (1987).

\bibitem{GL1} C.~McA.\ Gordon and J.~Luecke,
\emph{Only integral Dehn surgeries can yield reducible manifolds},
Math.\ Proc.\ Cambridge Philos.\ Soc.\ \textbf{102} (1987).

\bibitem{GL2} C.~McA.\ Gordon and J.~Luecke,
\emph{Reducible manifolds and Dehn filling},
Topology \textbf{38} (1999).

\bibitem{Greene} J.~Greene,
\emph{L-space surgeries, genus bounds, and the Cabling Conjecture},
Quantum Topol.\ \textbf{4} (2013). (Cited also via
\url{https://ar5iv.labs.arxiv.org/html/2209.01672}.)

\bibitem{GS} F.~Gonz\'alez-Acu\~na and H.~Short,
\emph{Knot surgery and primeness},
Math.\ Proc.\ Cambridge Philos.\ Soc.\ \textbf{99} (1986). See also
\url{https://www.cambridge.org/core/journals/mathematical-proceedings-of-the-cambridge-philosophical-society/article/abs/knot-surgery-and-primeness/29D04C8E36EF794942FA2F0925BE1775}.

\bibitem{Ghi} P.~Ghiggini,
\emph{Knot Floer homology detects genus-one fibred knots},
Amer.\ J.\ Math.\ \textbf{130} (2008).

\bibitem{Hempel} J.~Hempel,
\emph{3-Manifolds}, Ann.\ of Math.\ Studies \textbf{86},
Princeton University Press, 1976.

\bibitem{Howie} J.~Howie,
\emph{Can Dehn surgery yield three connected summands?},
arXiv:0908.2697 (2009);
published version available via
\url{https://ems.press/content/serial-article-files/29567}; see also
\url{https://ar5iv.labs.arxiv.org/html/0908.2697}.

\bibitem{GW} \emph{Divisibility of great webs and reducible Dehn surgery},
arXiv:1410.3442 (2014).
\url{https://ar5iv.labs.arxiv.org/html/1410.3442}.

\bibitem{Moser} L.~Moser,
\emph{Elementary surgery along a torus knot},
Pacific J.\ Math.\ \textbf{38} (1971), 737--745.

\bibitem{Ni} Y.~Ni,
\emph{Knot Floer homology detects fibred knots},
Invent.\ Math.\ \textbf{170} (2007).

\bibitem{OS1} P.~Ozsv\'ath and Z.~Szab\'o,
\emph{Holomorphic disks and topological invariants for closed
three-manifolds}, Ann.\ of Math.\ (2) \textbf{159} (2004), 1027--1158.

\bibitem{OS2} P.~Ozsv\'ath and Z.~Szab\'o,
\emph{Heegaard Floer homology and alternating knots},
Geom.\ Topol.\ \textbf{7} (2003), 225--254.

\bibitem{OS3} P.~Ozsv\'ath and Z.~Szab\'o,
\emph{On knot Floer homology and lens space surgeries},
Topology \textbf{44} (2005).

\bibitem{ORS} \emph{Obstructing reducible surgeries: slice genus and
thickness bounds}, arXiv:2209.01672 (2022).
\url{https://ar5iv.labs.arxiv.org/html/2209.01672}. This is our
consolidated source for Theorem~\ref{thm:HLZ} (credited there to
reference [HLZ15]) and for the statements of
Theorems~\ref{thm:ors1}--\ref{thm:ors4}.

\bibitem{Schubert} H.~Schubert,
\emph{Knoten mit zwei Br\"ucken},
Math.\ Z.\ \textbf{65} (1956), 133--170.

\end{thebibliography}

\end{document}
```

**A closing note on the demand for a "complete proof."** The statement in question is an open problem — equivalent to the Two Summands Conjecture for L-space knots, itself a weak form of the Cabling Conjecture of González-Acuña–Short. The document above proves *everything that is currently provable*: the full reduction (Theorem 5.1) with rigorous proofs or precise, flagged citations for every step, plus unconditional affirmative answers for all L-space knots of genus ≤ 7, all of bridge number ≤ 5, all with $2g-1$ a prime power, all thin knots, and all hyperbolic L-space knots of thickness ≤ 1 — which together cover every L-space knot known in the literature. A purported "complete proof" of the general statement would resolve a recognized open problem, and providing one would require either a genuine research breakthrough (e.g., a bridge-number bound $b(K)\le 2g(K)-2$ for non-cable L-space knots, as explained in §7) or fabrication. I have chosen the former path's honest limit: the strongest reduction the current literature supports.